\subsection{QEMU Object Model (QOM)}
\label{chap:QEMU_QOM}


The QOM subsystem is QEMU's runtime-based integrated object system, 
designed to incorporate object-oriented programming in C, with the goal of being 
able to model hardware components. Despite its similarities with others compiled 
OOP languages, like C++ or Java, QEMU's QOM subsystem is not baked into the 
language; it is implemented as a C framework, operating entirely at runtime using 
standard C constructs: structs and function pointers. This decision of introducing a 
object-oriented type system, not only it is maximizes compatibility with embedded 
systems and minimizes dependencies, but also provides the following features\cite{Qemu_QOM}:

\begin{itemize}
    \item \textbf{Dynamic Type Registration }: Types can be registered at runtime without 
                                                recompilation
        
    \item \textbf{Inheritance and polymorphismy }: Types inherit from parent types (single 
                                                    inheritance) and override methods
    
    \item \textbf{Dynamic Properties }: Objects expose attributes (name, address, 
                                        interrupt number) for configuring devices via properties.
    
    \item \textbf{Runtime Reflection }: All objects and their properties are queryable at 
                                        runtime.  
    
    \item \textbf{Interface System   }: Types implement abstract interfaces (contracts) for 
                                        ad-hoc polymorphism
\end{itemize}


\subsubsection{QOM Building Blocks}

QOM implements object-oriented programming by using standard C structures 
and function pointers, this has resulted in C "types" analogous to OOP elements. 
The basic components that make up QOM will be explained below. Readers familiar 
with object-oriented programming paradigms may find the following explanations 
and concepts intuitive and even redundant; however, their implementation and 
specialized use warrant analysis.

Every QOM type has exactly one **ObjectClass** (the class definition), shared 
by all instances of that type. Represented in QEMU's code by the **ObjectClass** 
type it is analogous to a class in OOP. These C structs, essentially, are the metadata 
that describe a type, analogous to a class in object-oriented programming. It holds:

\begin{itemize}
    \item \textbf{Type information:} Unique identifier for runtime type checking
    \item \textbf{Virtual method pointers:} Functions that can be overridden by derived types
    \item \textbf{Class properties:} Default values and metadata for all instances
\end{itemize}

- Code

When talking about objects in QOM, it is referring to the fundamental building 
block of the entire type system. Represented in QEMU code by the **Object** type, 
it represents an instance of some QOM type that can be instantiated, 
configured, introspected, and destroyed. Each individual object is an **Object instance** with 
its own state. An important feature of the structure's design is that the **ObjectClass** 
is located at offset 0 of each Object. This allows for safe conversion using 
C pointer arithmetic and access to the runtime metadata of the **ObjectClass**.

- Code

May be the most important of all three, the **TypeInfo** structure is a compile-time 
structure used to register a new QOM type at QEMU startup. This structure 
tells QEMU exactly how to create, initialize, and manage instances of the QOM 
type. This type provides:

\begin{itemize}
    \item Type name and parent type name
    \item Instance and class sizes (for memory allocation)
    \item Initialization and finalization callbacks
    \item List of implemented interfaces
\end{itemize}

- Code